% ----------------------------------------------------------------------------
%                               METODOLOGIA
% ----------------------------------------------------------------------------

\chapter{TECNOLOGIAS}
\label{chap:tecnologias}

Nesta seção, serão descritas as tecnologias necessárias para o desenvolvimento do Painel Parkinson em versão \textit{desktop}.

\section{\textit{Front-end}}
\label{sec:Front end}

O \textit{front-end} consiste na interação do usuário com o sistema. O \textit{front-end} tem como maior objetivo a experiência do usuário, manipulação das informações e a criação de identidade visual, levando em consideração \textit{User Experience} (UX) e suas boas práticas.

\subsection{Javascript}
\label{sec:Javascript}

JavaScript é uma linguagem de programação interpretada, baseada em objetos e amplamente usadas para páginas web. O JavaScript segue o padrão ECMAScript, que foi incorporado em todos os navegadores modernos.

\subsection{VueJS}
\label{sec:VueJS}

\textit{Framework front-end} que traz como principal vantagem ser \textit{Open Source}. Possui uma renderização declarativa que compreende a sintaxe de um modelo com saída HTML com base no JavaScript. Além disso, rastreando as mudanças de estados do JavaScript com atualizações do DOM de forma eficiente \cite{vuejs_intro}.

\subsection{Electron.js}
\label{sec:ElectronJS}

O Electron.js é um \textit{framework} capaz de criar aplicações com JavaScript, HTML e CSS. Com isso, permitindo uma base de código para o JavaScript, além disso, funciona em multiplataforma, sendo elas: Windows, macOS e Linux.
Possui uma grande confiabilidade, interoperabilidade, versatilidade e ubiquidade.

\subsection{Painel Parkinson}
\label{sec:PAINEL}

O módulo faz parte do Painel Parkinson. Um sistema web desenvolvido em VueJS que utiliza o SIDABI como banco de dados. O painel Parkinson foi desenvolvido para apoiar profissionais da saúde na análise de diagnóticos e apoio à decisões em pacientes  com a doença de Parkinson.

\section{\textit{Back-end}}
\label{sec:Back end}

O \textit{Back-end} é uma parte da aplicação que tem como principais responsabilidades, a execução lógicas mais complexas e o armazenamento dos dados para a aplicação.
\textit{Back-end}\cite{alura_backend}.

Sua função é de suma importância para a integração do \textit{front-end} com o banco de dados. Com isso, garantindo segurança das informações, performance e modularização do sistema.
As tecnologias no desenvolvimento \textit{back-end} incluem a linguagem Python com o \textit{framework} Flask e o PostgreSQL com ligação direta ao SIDABI.

\subsection{Python}
\label{sec:Python}
O Python possui uma sintaxe clara que favorece a legibilidade do código-fonte, tornando a linguagem mais produtiva \cite{borges2014python}. Além disso, é um \textit{software} de código aberto com licença compatível com a GPL \cite{borges2014python}.

Incluindo uma diversidade de estruturas de alto nível e inúmeros \textit{frameworks} que podem ser adicionados.

Além de ser uma das linguagens mais utilizadas pelos desenvolvedores que permite automatizar tarefas com o PostgreSQL.

Essa produtividade será essencial para a produção do Painel Parkinson em forma \textit{desktop} no processo de automação das informações e as regras de negócios.

\subsection{Flask}
\label{sec:Flask}

O Flask é um \textit{framework} leve para aplicações web WSGI. Ele facilita e agiliza o desenvolvimento sendo poderoso por ser um \textit{framework} expansível que ajuda desde o início do desenvolvimento e um núcleo poderoso para a instalação do que for necessário para a aplicação do Painel Parkinson\cite{grinberg2018flask}.

\subsection{PostgreSQL}
\label{sec:PostgreSQL}

Segundo a documentação oficial, "O PostgreSQL é um poderoso sistema de banco de dados objeto-relacional de código aberto que utiliza e estende a linguagem SQL, combinada com diversos recursos que permitem armazenar e dimensionar com segurança as cargas de trabalho de dados mais complexas" \cite{postgresql_about}.

O PostgreSQL é um \textit{software} gratuito e de código aberto e conta com soluções eficientes, integridade dos dados e alto desempenho.